\section{Discussion}
\label{sec:discussion}

The transition in $X_9$ exhibits a compact coupling between
two conifold sectors and an interacting $E_2$ sector. Its
charge relation constrains the invariant Higgs coordinates,
its integral neutral lattice fixes the surviving vector
fields, and its mirror periods determine their genus-zero
quantum couplings. The four-hypermultiplet thresholds reproduce the
singular monodromy, while the eight-period comparison fixes the
regular scalar metric and electromagnetic couplings of the neutral
sector. Absence of neutral logarithmic thresholds alone would not
establish that comparison.

The two limits explain the scope of this effective description.
With bounded local K\"ahler scales, sending the compact volume to
infinity removes the shared flavor gauging. Sending the auxiliary
radius to infinity at fixed five-dimensional Planck scale retains
finite bulk gauging but closes the Kaluza--Klein gap. Neither limit
produces an autonomous four-particle description retaining the
original cooperative constraint.

Three extensions would materially strengthen the physical
description. First, an explicit finite-Planck hypermultiplet
geometry would test how the rigid invariant-coordinate
description is embedded in quaternionic-K\"ahler geometry.
Second, a semistable heterotic degeneration with bundle data
would determine the microscopic realization of the shared
gauging, including the masslessness of its $U(1)$. The K3
slice and topological counts alone do not provide this.
Third, a stability condition along the quantum continuation
would distinguish which magnetic sheaves remain physical
particles and which can decay. The isolated numerical
Donaldson--Thomas contributions and partial wall exclusions
in \cite{Wuebben:2026magnetic} do not determine that spectrum.

Geometric existence is governed by the exact mixed smoothing
criterion, including the three discriminant conditions on the
$SU(3)$ component of the $E_3$ branch
\cite{Wuebben:2026smoothing}. Selected-base integrability under its
projection hypothesis removes a geometric obstruction; it does not
construct the compact hypermultiplet action. Questions about selected
profiles with zero $E_3$ projection and unrestricted nonreduced
deformation spaces remain outside that theorem. On the quantum side, continuation beyond the
marked branch and higher-genus transition formulas are natural
next problems. The results proved here require neither of
these extensions: they establish one actual mixed compact
transition, the obstruction to retaining its flavor gauging
in the specified rigid limits, and the complete genus-zero
vector geometry across its circle-reduced transition.
